\section*{第 \ref{ch:b3:membrane-traffic} 章——膜泡运输与蛋白质分选}

\begin{solution}{exo:b3:membrane-traffic:1}
氨基端疏水区段：\omterm{def:b3:membrane-traffic:signals}{信号肽}，进入
内质网（默认情况下继而运至细胞表面）。PKKKRKV：
\omterm{def:b3:membrane-traffic:signals}{核定位信号}，经核孔进入细胞核。
氨基端两亲螺旋：线粒体\omterm{def:b3:membrane-traffic:signals}{前导序列}，进入
线粒体基质。KDEL：从高尔基体回收至
内质网。\omterm{prop:b3:membrane-traffic:golgi}{甘露糖-6-磷酸}：从反面高尔基体运往\omterm{def:b3:membrane-traffic:endocytosis}{内体}
和\omterm{def:b3:membrane-traffic:endocytosis}{溶酶体}。
\end{solution}

\begin{solution}{exo:b3:membrane-traffic:2}
（1）无膜时，翻译所得肽链比分泌的
蛋白质长约二十个残基：
分泌过程中该片段被切除。（2）翻译时若存在膜，产物
大小与成熟蛋白相同，且不被蛋白酶降解：它已进入
囊泡，并在其中失去该片段。（3）翻译后才加入的膜
两者均不发生：进入过程与合成相耦联，
即共翻译——，该片段（\omterm{def:b3:membrane-traffic:signals}{信号肽}）仅作用于
新生肽链。
\end{solution}

\begin{solution}{exo:b3:membrane-traffic:3}
内质网至高尔基体：\omterm{def:b3:membrane-traffic:coats}{COPII}，顺向运输；高尔基体至内质网：\omterm{def:b3:membrane-traffic:coats}{COPI}，
逆向运输。反面高尔基体至\omterm{def:b3:membrane-traffic:endocytosis}{内体}：\omterm{def:b3:membrane-traffic:coats}{网格蛋白}及其衔接蛋白，
顺向运往\omterm{def:b3:membrane-traffic:endocytosis}{溶酶体}。质膜至\omterm{def:b3:membrane-traffic:endocytosis}{内体}：\omterm{def:b3:membrane-traffic:coats}{网格蛋白}，
内向运输（胞吞作用）。
\end{solution}

\begin{solution}{exo:b3:membrane-traffic:4}
糖类只在内质网和高尔基体的腔内加入。每一囊泡的
腔面在出芽与融合过程中始终朝向腔内，
囊泡与质\omterm{def:b3:membrane-traffic:fusion}{膜融合}时，其腔面成为
细胞外侧，取向从不颠倒，因此
原先在内侧被糖基化的部分最终位于外侧。
\end{solution}

\begin{solution}{exo:b3:membrane-traffic:5}
$t^{*} = \ln(0.05/0.2)/(0.05 - 0.2) = \ln 4/0.15 = \qty{9.2}{min}$;
$B(t^{*}) = (0.2/(-0.15))(e^{-1.85} - e^{-0.46}) = 0.63$。在 \qty{60}{min} 时
$A \approx 0$，$B = 1.33\,(e^{-3} - e^{-12}) = 0.066$，因此 $C = 0.93$。
\end{solution}

\begin{solution}{exo:b3:membrane-traffic:6}
\omterm{def:b3:membrane-traffic:signals}{信号肽}使氨基端位于腔内；每个
疏水区段交替地作为终止转移或起始转移序列，
因此共有四个跨膜螺旋（约 30--52、
90--112、150--172、210--232 位残基）。各片段分别为：氨基端位于腔面；
52--90 位于胞质侧；112--150 位于腔面；172--210 位于胞质侧；羧基
端位于腔面。糖基化仅发生在腔面片段上（即
氨基端、112--150 片段和羧基端），这些片段在
表面变为胞外侧。
\end{solution}

\begin{solution}{exo:b3:membrane-traffic:7}
$J_{\max} = \num{30000}\times 0.25\times 0.05/0.30 = 1250$ 每分钟；
表面比例 $k_{r}/(k_{i} + k_{r}) = 0.05/0.30 = 0.17$。$R$ 加倍
则 $J$ 恰好加倍；回流速率是瓶颈（多数受体
位于细胞内），若将 $k_{r}$ 加倍，得 $\num{30000}\times 0.25
\times 0.1/0.35 \approx 2140$，即为 $1.7$ 倍。
\end{solution}

\begin{solution}{exo:b3:membrane-traffic:8}
（a）无受体：带标签的水解酶走默认途径，
被分泌到细胞外；\omterm{def:b3:membrane-traffic:endocytosis}{溶酶体}缺乏酶，充满未被消化的
物质——即一种贮积病。（b）无磷酸转移酶：
表型相同，但病因不同（I 细胞病），酶
未被标记而被分泌。（c）\omterm{def:b3:membrane-traffic:endocytosis}{内体}被中和：受体无法
释放其货物，因而不能回收利用，水解酶被错误投送
或被分泌；LDL 及其他受体也停止循环。
\end{solution}

\begin{solution}{exo:b3:membrane-traffic:9}
该酪氨酸属于\omterm{def:b3:membrane-traffic:coats}{网格蛋白}衔接蛋白所识别的尾部\omterm{def:b3:bioinformatics:motif}{基序}；
缺少它，受体便不能聚集到被衣小窝中。
受体与 LDL 结合正常，但均匀分布于
表面，不进入小窝，配体留在细胞外——
这是分选障碍，而非结合障碍。
\end{solution}

\begin{solution}{exo:b3:membrane-traffic:10}
若 $k_{1} = k_{2} = k$，试取 $B = kt\,e^{-kt}$：$\mathrm{d}B/\mathrm{d}t
= k e^{-kt} - k^{2}t e^{-kt} = kA - kB$，符合要求，且 $B(0) = 0$。
峰值位于 $1 - kt = 0$ 处：$t^{*} = 1/k$，$B = e^{-1} = 0.37$。一般地
$C(t) = 1 - (k_{2}e^{-k_{1}t} - k_{1}e^{-k_{2}t})/(k_{2} - k_{1})$,
表达式关于 $k_{1}$ 与 $k_{2}$ 互换具有对称性：仅凭
颗粒的曲线无法判断两个区室中哪一个是快速的。
\end{solution}

\begin{solution}{exo:b3:membrane-traffic:11}
每个囊泡的面积为 $4\pi(50\,\text{nm})^{2} = \qty{0.031}{\micro
m^2}$；每分钟一千个即为 \qty{31}{\micro m^2/min}，因此整个
表面所需时间为 $3000/31 \approx \qty{95}{min}$。细胞不会缩小，
因为膜通过胞吐作用（回收囊泡）以相同
速率返回——即流量的平衡。每个囊泡的体积为 $5\times 10^{-19}$ L；
每分钟一千个，持续一小时，共 $3\times 10^{-14}$ L，约
相当于一个 \qty{2}{pL} 细胞每小时体积的 \qty{1.5}{\%}。
\end{solution}

\begin{solution}{exo:b3:membrane-traffic:12}
肉毒杆菌毒素作用于运动终末，乙酰胆碱不再释放，
肌肉得不到指令而松弛——即弛缓性
麻痹。破伤风毒素沿运动轴突上行，再转运至
抑制性中间神经元，它们通常抑制运动神经元；
其 \omterm{def:b3:membrane-traffic:fusion}{SNARE} 被切割后，甘氨酸或 GABA 不再释放，
运动神经元不受制约地放电，所有肌肉收缩——即痉挛性麻痹。
少量（数纳克）肉毒杆菌毒素注入过度活跃的肌肉后，
可使其局部失能数月（直至新的 \omterm{def:b3:membrane-traffic:fusion}{SNARE} 合成、
终末恢复）：可用于治疗肌张力障碍、痉挛、斜视与皱纹。
\end{solution}

\begin{solution}{pb:b3:membrane-traffic:1}
\textbf{1.} $10^{7}\times \qty{25000}{Da}\times 1.66\times 10^{-24}$ 克
$= \qty{0.42}{pg}$ 每分钟，\qty{600}{pg} 每天——为其自身
蛋白质含量的两倍。
\textbf{2.} 囊泡体积 $\tfrac{4}{3}\pi\,35^{3} = 1.8\times 10^{5}$
nm$^{3}$，其 \qty{30}{\%} 为 $5.4\times 10^{4}$；一个酶分子
$\tfrac{4}{3}\pi\,2^{3} = 34$ nm$^{3}$：每个囊泡约含 \num{1600} 个
分子；$10^{7}/1600 \approx 6200$ 个囊泡每分钟。
\textbf{3.} 每个囊泡面积 $4\pi\,35^{2} = \qty{0.0154}{\micro m^2}$，每分钟输出 \qty{95}{\micro
m^2} 的膜，若无回收，内质网的 \qty{6000}{\micro m^2} 将在
约一小时内耗尽。
\textbf{4.} $\qty{95}{\micro m^2} = 9.5\times 10^{7}$ nm$^{2}$，两层小叶
各以 \qty{0.7}{nm^2} 计：$2.7\times 10^{8}$ 个脂质分子每分钟。细胞
借 \omterm{def:b3:membrane-traffic:coats}{COPI} 囊泡与回收利用途径收回膜，并合成脂质，
以补充随颗粒永久离开的脂质。
\textbf{5.} 颗粒体积 $5.2\times 10^{8}$ nm$^{3}$，其 \qty{30}{\%} 为
$1.6\times 10^{8}$：每个颗粒 $4.7\times 10^{6}$ 个分子；$6\times
10^{8}$ 个分子每小时约可充满 $130$ 个颗粒。
\textbf{6.} 每个颗粒增加 $\pi d^{2} = \qty{3.1}{\micro m^2}$；
$300$ 个颗粒共增加 \qty{940}{\micro m^2}，使 \qty{50}{\micro m^2} 增至原来的
二十倍。膜须于数分钟内经补偿性胞吞作用收回。
\textbf{7.} $t^{*} = \ln 2/0.05 = \qty{13.9}{min}$，$B = 0.50$。
\textbf{8.} $A = e^{-3} = 0.05$；$B = 2(e^{-1.5} - e^{-3}) = 0.35$；
$C = 0.60$。
\textbf{9.} \qty{10}{min} 与 \qty{20}{min}；平均总时长 \qty{30}{min}。
\textbf{10.} $t^{*} = \ln(0.1)/(-0.09) = \qty{25.6}{min}$；$B(t^{*}) =
(0.1/(-0.09))(e^{-2.56} - e^{-0.256}) = 0.77$：高尔基体膨大，
充满了货物的四分之三。
\textbf{11.} 在峰值处，$k_{1}e^{-k_{1}t^{*}} = k_{2}e^{-k_{2}t^{*}}$；
代入 $B$ 得 $B_{\max} = (k_{1}/k_{2})^{\,k_{2}/(k_{2} -
k_{1})}$：指数为 $k_{2}/(k_{2} - k_{1})$。若 $k_{1} > k_{2}$ 时，
则指数为负，底数大于 $1$；若 $k_{1} < k_{2}$ 时，
则底数小于 $1$，指数为正——两种情形下，该值
均小于 $1$（验算：此处 $2^{-1} = 0.5$）。
\textbf{12.} 当 $k_{1} \gg k_{2}$ 时，$B(t) \to e^{-k_{2}t}$：标记
即刻进入高尔基体，并以 $k_{2}$ 速率离开——内质网一步
不可见，高尔基体的曲线为简单衰减。
\textbf{13.} $\theta = 2/(2+2) = 0.5$；$J = \num{50000}\times 0.2\times
0.1\times 0.5/(0.1 + 0.1) = 2500$ 个颗粒每分钟，每小时
\num{150000} 个。
\textbf{14.} $1.5\times 10^{5}\times 1500 = 2.3\times 10^{8}$ 个胆固醇
分子每小时；$5\times 10^{9}/2.3\times 10^{8} \approx 22$ 小时。
\textbf{15.} 表面占比 $k_{r}/(k_{r} + k_{i}\theta) = 0.5$；一次
往返循环耗时 $1/(k_{i}\theta) + 1/k_{r} = 10 + 10 = \qty{20}{min}$：
二十小时内约 $60$ 次往返。
\textbf{16.} $J = 1250$ 每分钟。未饱和时 $J \propto R\,[L]$，
故受体减半时，血浆 LDL 须加倍，
方能维持相同的清除速率。
\textbf{17.} $R$ 恢复至 \num{50000}：清除能力恢复，血浆 LDL
回落至正常水平——即他汀类药物的作用。
\textbf{18.} 若无受体，LDL 仅经缓慢的非特异途径清除，
因而循环数日而非数小时，
被氧化，并为巨噬细胞的清道夫受体所摄取，
这些巨噬细胞遂成为动脉粥样硬化斑块中的泡沫细胞；而他汀类药物
通过诱导受体发挥作用，此时已无受体可供诱导。
\textbf{19.} $V = \tfrac{4}{3}\pi(0.25)^{3} = \qty{0.065}{\micro m^3} =
6.5\times 10^{-17}$ 升。在 pH $7.2$ 时：$6.3\times 10^{-8}\times 6.5\times
10^{-17} = 4\times 10^{-24}$ 摩尔——约 $2$ 个游离质子。在 pH $4.7$ 时：
$2\times 10^{-5}\times 6.5\times 10^{-17} = 1.3\times 10^{-21}$ 摩尔，
约 $780$ 个。
\textbf{20.} $50\times 780 \approx \num{39000}$ 个质子；$50$ 个泵
各以每秒 $200$ 个的速度工作，每秒共转运 \num{10000} 个，约需 \qty{4}{s}。
\textbf{21.} 泵入正电荷使管腔带正电，
数千个电荷即足以使泵停止工作；氯离子经
通道进入（或阳离子外流）以中和电荷。
\textbf{22.} 在中性 pH 下失活的酶，若\omterm{def:b3:membrane-traffic:endocytosis}{溶酶体}泄漏，
或被错误分选至细胞溶胶或细胞表面，并不造成危害：
酸性条件将消化限定于为其而建的区室之内。
\textbf{23.} 二者都须带有 pH 感受器——即组氨酸，其
在 pH~6 附近的质子化会改变结合表面：LDL 受体
质子化时，将一个结构域折叠覆盖于自身的结合位点上，
\omterm{prop:b3:membrane-traffic:golgi}{甘露糖-6-磷酸}受体同样会失去亲和力。
\textbf{24.} $10^{7.2 - 4.7} = 10^{2.5} \approx 300$ 倍的浓缩；
被截留的碱消耗质子，抬高\omterm{def:b3:membrane-traffic:endocytosis}{溶酶体}的 pH，并
使水解酶失活——这正是氯喹在食物液泡中杀灭疟原虫的机理，
也是它在
实验中阻断\omterm{def:b3:membrane-traffic:endocytosis}{自噬}的方式。
\textbf{25.} 高尔基体的峰值出现于 \qty{14}{min}；约 \num{6200} 个囊泡
每分钟离开内质网；每小时被摄取的 LDL 颗粒
共 \num{150000} 个。
\end{solution}
